\section{Design}
\label{sec:design}

\subsection{Values from tables}

A specification file lists the CSV tables to expose, the columns that identify a row and a format
for each column (Listing~\ref{lst:spec}). \texttt{provenum values} reads the tables and writes one
LaTeX file that defines a value for every cell. The manuscript loads that file and, where it
would have written \texttt{+6.8}, writes
\begin{center}
\texttt{\textbackslash val\{comparisons/vina/top1/box\_padding\_2/delta\_pts\}}.
\end{center}
The address is the table name, the key values in order, and the
column.

\begin{lstlisting}[caption={Part of a specification. Proportions default to percentages with one
decimal; other columns are formatted by name.},label={lst:spec},float=t]
[[table]]
file = "comparisons.csv"
keys = ["engine", "metric", "arm"]
default = "pct1"
formats = { delta_pts = "+f1", p_holm = "p", gained = "int" }
\end{lstlisting}

Formatting is part of the specification rather than the text because the same cell can be
reported in more than one way. The formats cover the forms we met in practice: percentages
(\texttt{pct0}--\texttt{pct2}), fixed decimals (\texttt{f0}--\texttt{f4}), significant figures,
integers, p-values (two significant figures, scientific notation below 0.001), a \texttt{+} prefix
that always shows the sign, and a \texttt{\textasciitilde} prefix for ratios and other magnitudes
quoted approximately. Values are wrapped in \texttt{\textbackslash ensuremath} so they work in
text and in mathematics. An address that does not exist prints a bold ``??'' in the PDF and a
LaTeX warning, so a broken reference is visible at the next build.

\subsection{Lint}

\texttt{provenum lint} lists every number in the manuscript that is not produced by
\texttt{\textbackslash val\{\}}. Integers below 10 and plausible years are skipped as too
ambiguous to trace. With \texttt{--strict} it exits with an error if any hand-typed number
remains, which lets a continuous-integration job keep a migrated paper migrated. A GitHub Action
regenerates the values file from the tables and fails if it differs from the committed one, so
a table changed without rebuilding the values is caught as well.

\subsection{Migration: model proposes, code verifies}
\label{sec:anchor}

Migration needs, for every hand-typed number, the cell it reports. \texttt{provenum anchor}
extracts each number with its sentence and paragraph, marks the number in the sentence, and asks
a language model for a cell address. The model sees a catalogue of the addressable values: for
each table, the address pattern, the allowed key values and the columns, but not the cell values
themselves. It answers in a fixed JSON schema with an address or \texttt{null}.

The answer is checked by code. \texttt{judge} looks up the proposed cell, formats its value the
way the text writes the number (same number of decimals, sign and notation), and compares. There
are four outcomes: the value reproduces the text (\emph{traced}); the address exists but the
value differs (\emph{mismatch}); the address does not exist (\emph{bad path}); or the model
declined (\emph{abstained}). \texttt{provenum migrate} replaces only traced numbers. A mismatch is
reported for the author to read: either the model chose the wrong cell or the text is wrong, and
the program cannot tell which.

We call this blind mode, because the model chooses without seeing values. A second mode,
shortlist mode, is for migration only. The program first finds every cell whose formatted value
equals the written number, and the model chooses among those cells or answers ``none''. This is a
much easier task, but it cannot find errors: a wrong number has no matching cell, so it is never
offered. Shortlist mode is therefore used to convert numbers that are already right, and blind
mode to look for numbers that are wrong.

The model is never asked whether a number is correct. Its only output is an address, and an
address is either confirmed by the value or not. This bounds the damage a model can do. In blind
mode a wrong address cannot pass the value check unless the wrong cell happens to hold the same
rounded value. That case does occur (Section~\ref{sec:results}), and it is not harmless: the text
renders correctly today but would follow the wrong cell after a rerun. We return to it in
Section~\ref{sec:discussion}.

\subsection{Audit without a model}

\texttt{provenum audit} checks numbers against the tables without changing the text and without a
model. It scores each cell that could match a number by the overlap between the sentence's words
and the cell's key values and column name, after a project vocabulary maps phrases such as
``2~\AA{} box'' to key values such as \texttt{box\_padding\_2}. We include it as a baseline for
anchoring.

\subsection{Local models}

The model interface uses the Ollama HTTP API through the Python standard library, so \pv{} has no
runtime dependencies and the manuscript need not leave the author's machine. Requests use
temperature 0 and a JSON schema for the reply. For Qwen3 models the reasoning phase is turned off,
since the other models answer directly and the task is a single choice.
